\Askhsh{31534_SOLUTION}{1}
\begin{ylh}{1}
    \item [ΛΥΣΗ]
\end{ylh}
\begin{alist}
\item
Από τα δεδομένα του σχήματος προκύπτει ότι $f'(x)=\alpha x^2+\beta x+\gamma$, $\alpha<0$ και $f'(0)=0$, οπότε $\gamma=0$. Η συνάρτηση $f'$ παρουσιάζει μέγιστο για $x=2$, οπότε $f''(2)=0$. Είναι: $f'(x)=2\alpha x+\beta$ και
\[ f''(2)=0 \Leftrightarrow 4\alpha+\beta=0, (1) \]
Επιπλέον, $f'(2)=2 \Rightarrow 4\alpha+2\beta=2, (2)$.
Αν από την (2) αφαιρέσουμε την (1) παίρνουμε $\beta=2$ και με αντικατάσταση στην (1) βρίσκουμε $\alpha=-\frac{1}{2}$.
Ετσι έχουμε $f'(x)=-\frac{1}{2}x^2+2x$. \monades{X}

\item
Από το ερώτημα α) συμπεραίνουμε ότι η $f(x)=\frac{1}{3}x^3+x^2+x+c$, $c \in \mathbb{R}$ και $f(0)=1$, οπότε $c=1$. Άρα $f(x)=-\frac{1}{6}x^3+x^2+c$. \monades{X}

\item
\begin{rlist}
\item
Αρκεί να αποδείξουμε ότι για κάθε $x>0$ ισχύει
\[ x^2+x+1-\eta\mu x > -\frac{1}{6}x^3+x^2+1 \]
που γράφεται
\[ \frac{1}{6}x^3+x-\eta\mu x > 0 \]
Θεωρούμε τη συνάρτηση $h(x)=\frac{1}{6}x^3+x-\eta\mu x$, $x \ge 0$.
Η συνάρτηση είναι παραγωγίσιμη στο πεδίο ορισμού της με $h'(x)=\frac{1}{2}x^2+1-\sigma\upsilon\nu x$ και $h'(x) \ge 0$ με την ισότητα να ισχύει μόνο για $x=0$. Άρα η $h$ είναι γνησίως αύξουσα, οπότε για κάθε $x>0$ έχουμε:
\[ h(x)>h(0) \Rightarrow \frac{1}{6}x^3+x-\eta\mu x > 0 \]
που είναι το ζητούμενο.

Εναλλακτικά θα μπορούσαμε να γράψουμε:
Για οποιοδήποτε θετικό αριθμό $x$ έχουμε:
\[ |\eta\mu x| < x \Leftrightarrow -x < \eta\mu x < x \]
και επειδή με $x>0$ ισχύει: $\frac{1}{6}x^3>0$. Άρα, $\frac{1}{6}x^3+x-\eta\mu x > 0$. \monades{X}

\item
Από το προηγούμενο ερώτημα συμπεραίνουμε ότι το εμβαδόν του χωρίου που ορίζεται από τις $C_f$, $C_g$ και τις ευθείες $x=0$ και $x=\pi$ είναι ίσο με $\int_0^\pi h(x)dx$. Είναι:
\[ \int_0^\pi h(x)dx=\int_0^\pi \left(\frac{1}{6}x^3+x-\eta\mu x\right)dx=\left[\frac{x^4}{24}+\frac{x^2}{2}+\sigma\upsilon\nu x\right]_0^\pi \]
\[ =\frac{\pi^4}{24}+\frac{\pi^2}{2}-1-1=\frac{\pi^4+12\pi^2-48}{24} \] \monades{X}
\end{rlist}
\end{alist}